\subsubsection{Оптимізація алгоритму: кешування ваг та багатопотокова обробка}

Фундаментальним недоліком прямого методу є багаторазове і надлишкове обчислення вагових коефіцієнтів у глибоко вкладених циклах. Оскільки ваги Чебишева залежать виключно від просторових координат вузлів сітки і не залежать від значень кольору самих пікселів, їх можна обчислити заздалегідь. Ця властивість стала основою для розробки оптимізованого алгоритму (\texttt{processBarycentricZoom}).

Першим етапом оптимізації є попереднє обчислення та збереження ваг. У програмному коді реалізовано спеціалізовану структуру даних \texttt{WeightCache}, яка зберігає масив вагових коефіцієнтів, їхню загальну суму (знаменник барицентричної формули), а також логічний прапорець \texttt{hit} та індекс \texttt{exact\_index}. Останні використовуються для миттєвої ідентифікації точного збігу координат цільової точки з вузлом вихідної сітки, що дозволяє безпечно уникати ділення на нуль без складних розгалужень у внутрішніх циклах обробки пікселів. Для максимального прискорення ініціалізації, розрахунок кеш-масивів для осей X та Y (\texttt{W\_X} та \texttt{W\_Y}) виконується паралельно у двох незалежних потоках (\texttt{std::thread}).

Другим етапом оптимізації стало усунення чотирьох вкладених циклів, які створювали надмірне обчислювальне навантаження. Замість них реалізовано двопрохідну схему обчислень із виділенням проміжного буфера пам'яті, що розділяє роботу алгоритму на два кроки:
\begin{enumerate}
	\item \textbf{Перший прохід (\texttt{pass1}):} Алгоритм ітерується по цільовій висоті $Y$ та вихідній ширині $i$, обчислюючи часткові барицентричні суми. Результати записуються у проміжний двовимірний масив \texttt{tempImg} розмірністю $H_2 \times (n+1)$. Для уникнення втрати точності на етапі проміжних обчислень, колірні компоненти RGB у \texttt{tempImg} зберігаються у форматі чисел з плаваючою комою (\texttt{float}).
	\item \textbf{Другий прохід (\texttt{pass2}):} Використовуючи попередньо підготовлені високоточні дані з \texttt{tempImg}, алгоритм виконує фінальне обчислення для кожної цільової координати $X$. Отримані значення нормалізуються, обмежуються стандартним колірним діапазоном $[0, 255]$ за допомогою функції \texttt{std::clamp} та записуються безпосередньо у результуючий об'єкт \texttt{QImage}.
\end{enumerate}

Ключову роль у забезпеченні швидкодії відіграє паралельна обробка даних на рівні центрального процесора. За допомогою визначення кількості доступних апаратних ядер (\texttt{std::thread::hardware\_concurrency}), виконання обох проходів розпаралелюється. Матриця цільового зображення розбивається на горизонтальні сегменти (chunks) вздовж осі Y. Кожен обчислювальний потік отримує свій незалежний блок рядків для обробки. Такий підхід гарантує, що потоки здійснюють запис у неперетинні області пам'яті, що повністю усуває необхідність використання механізмів блокування (наприклад, \texttt{std::mutex}) та запобігає виникненню стану гонитви (race condition).

Реалізація структурованого кешування у поєднанні з багатопотоковою двопрохідною архітектурою кардинально знизила обчислювальне навантаження. Замість екстремальної початкової складності $\mathcal{O}(W_2 \cdot H_2 \cdot W \cdot H)$, загальна кількість операцій тепер визначається сумою складностей двох послідовних проходів: $\mathcal{O}(H_2 \cdot W \cdot H)$ для першого етапу та $\mathcal{O}(H_2 \cdot W_2 \cdot W)$ для другого. Таким чином, загальна асимптотична складність оптимізованого алгоритму становить $\mathcal{O}(H_2 \cdot W \cdot (H + W_2))$. 

Для порівняння: при масштабуванні зображення розміром $500 \times 500$ пікселів у розмір $1000 \times 1000$, кількість ітерацій базового циклу зменшилася з 250 мільярдів (у прямому методі) до лише 750 мільйонів. Таке математичне скорочення кількості операцій більш ніж у 330 разів дозволило ефективно утилізувати апаратні ресурси багатоядерного процесора. Це перетворило алгоритм на практичне рішення, здатне обробляти зображення високої роздільної здатності за частки секунди без блокування графічного інтерфейсу користувача.